\documentclass[11pt]{article}
% Shared preamble for per-Python-file documentation (input via \input{...}).
\usepackage[margin=1in]{geometry}
\usepackage[T1]{fontenc}
\usepackage[utf8]{inputenc}
\usepackage{lmodern}
\usepackage{hyperref}
\usepackage{xcolor}
\usepackage{listings}
\usepackage{listingsutf8}
\lstset{
  basicstyle=\ttfamily\small,
  columns=fullflexible,
  breaklines=true,
  upquote=true,
  inputencoding=utf8,
  extendedchars=true,
  frame=single,
  framerule=0.2pt,
  tabsize=2
}


\title{\texttt{run\_d2co.py} (Q\&A)}
\author{}
\date{}

\begin{document}
\maketitle

\paragraph{Q: What is the purpose of \texttt{run\_d2co.py}?}
\textbf{A:} It trains and evaluates a Wide \& Deep model using a \textbf{GMM-based label transformation} conditioned on duration buckets. The script fits Gaussian Mixture Models per bucket to estimate two modes of play-time, maps labels into a bounded target space, trains with MSE, and maps predictions back to the original scale for evaluation.

\paragraph{Q: What is the key idea behind the GMM label transformation here?}
\textbf{A:} For each duration bucket, the script fits a 2-component Gaussian Mixture Model to play-time values, producing two means (interpreted as ``negative'' and ``positive'' modes). Labels are transformed using an exponential mapping (\lstinline|get_gmm_label|), and predictions are inverted back via \lstinline|get_real_value|.

\paragraph{Q: How are GMM means computed and stabilized?}
\textbf{A:} The function \lstinline|get_gmm_mean| iterates through buckets, fits a GMM for each, collects the sorted means, and then applies a frequency-weighted moving average smoothing (windowed rolling average) to reduce noise in sparse buckets.

\paragraph{Q: What model is trained in this script?}
\textbf{A:} It constructs \lstinline|WideAndDeep| (from \lstinline|model/wd.py|) and trains it with MSE against the transformed GMM label target.

\paragraph{Q: How does evaluation work?}
\textbf{A:} At test time, it uses the inverse transform to map the model's output back to a play-time value and then computes MAE/XAUC/KL (with dataset-specific MAE rescaling).

\paragraph{Q: Code example: fitting per-bucket GMM means (schematic)?}
\textbf{A:}
\begin{lstlisting}[language=Python]
gm = GaussianMixture(n_components=2, covariance_type="spherical", ...).fit(playTimeInBucket)
means = np.sort(gm.means_.T[0])
\end{lstlisting}

\paragraph{Q: Full source code?}
\textbf{A:}
\lstinputlisting[language=Python]{/workspace/run_d2co.py}

\end{document}
